\setcounter{section}{0}

\Section{Действительное число. 
Теорема о существовании и единственности точной верхней (нижней) 
грани ограниченного множества. Счетность \(\Q\), несчетность \(\R\)}

\Subsection{Действительное число}

\begin{definition}
    \textit{Рациональным числом} называют дробь вида \(\frac{p}{q}\), где \(p\in\Z,~q\in\N\). 
\end{definition}

\begin{definition}
    \textit{Сечением множества рациональных чисел \(\Q\)} называют
    разбиение этого множества на два непустых множества \(A_{*}\) и \(A^{*}\) таких, что:
    \begin{itemize}
        \item \(A_{*}\cup A^{*}=\R\) 
        \item \(A_{*}\cap A^{*}=\emptyset\)
        \item \(\forall\, x\in A_{*},~\forall\,y\in A^{*}\longmapsto x < y\)
    \end{itemize}
    При этом множества \(A_{*}\) и \(A^{*}\) называют
    \textit{нижним} и \textit{верхним классами сечения} соответственно. 
\end{definition}

\begin{notation}
    Будем называть сечения, в которых существует наибольший элемент в \(A_{*}\), 
    но не существует наименьший элемент в \(A^{*}\), \textit{сечениями класса А}.

    \vspace{2mm}

    Если же, наоборот, отсутствует наибольший элемент в \(A_{*}\), но есть наименьший
     элемент в \(A^{*}\), то будем относить такие сечения к \textit{классу В}.

     \vspace{2mm}

    В случае, когда нет ни наибольшего элемента в \(A_{*}\), ни наименьшего в \(A^{*}\),
    будем относить такие сечения к \textit{классу С}.
\end{notation}

\begin{definition}
    \textit{Иррациональным числом} называется сечение класса \(C\). \\
    \textit{Действительным числом} называются сечения любого из трех классов.
\end{definition}

\begin{definition}
    \textit{Сечение множества действительных чисел \(\R\)} определяется 
    аналогично сечению \(\Q\).
\end{definition}

\begin{definition}
    Множество \(X\subset\R\) называют \textit{ограниченным сверху}, 
    если \(\exists\,\alpha\in\R\) такое, что для любого элемента \(x\in X\)
    выполняется \(x\leq\alpha\). \\
    При этом число \(\alpha\) называют \textit{верхней гранью} множества \(X\).
    Ограниченность снизу и нижняя грань определяются аналогично.
    
    \vspace{2mm} 

    Множество \(X\subset\R\) называют \textit{ограниченным}, 
    если \(\exists\,\alpha\in\R\) такое, что для любого элемента \(x\in X\)
    выполняется \(|x|\leq\alpha\). \\
    Можно также сказать, что если множество органичено сверху и снизу, то оно ограничено.
\end{definition}

\begin{definition}
    Наименьшую из верхних граней множества \(X\) называют \textit{точной верхней гранью}.
    \[[\alpha_{0} = sup X]\overset{def}{=}[\forall\,x\in X\longmapsto\alpha_{0}\geq x]
    ~\&~[\forall\,\alpha\in\R, \alpha<\alpha_{0}~\exists\,x_{\alpha}\in X:x_{\alpha}>\alpha]\]
    Аналогично определяется \textit{точная нижняя грань}.
    \[[\alpha_{0} = inf X]\overset{def}{=}[\forall\,x\in X\longmapsto\alpha_{0}\leq x]
    ~\&~[\forall\,\alpha\in\R, \alpha>\alpha_{0}~\exists\,x_{\alpha}\in X:x_{\alpha}<\alpha]\]
\end{definition} 

\begin{lemma}
    Для любых действительных чисел \(\alpha\) и \(\beta\) таких, что \(\alpha>\beta\), существует 
    \textbf{рациональное} число \(c\), для которого верно соотношение \(\alpha>c>\beta\).

    \vspace{2mm}

    \proofname~Выберем по аксиоме Архимеда натуральное число \(n > \frac{1}{\alpha-\beta}\). \\
    Тогда имеем: \[na+1<nb\] 
    Теперь пусть целое число \(m = [n\alpha]+1\). Следовательно, справедливо следующее:
    \[m-1\leq n\alpha < m\]
    \[n\alpha < m \leq n\alpha + 1\]
    Теперь докажем, что число \(\frac{m}{n}\) лежит между \(\alpha\) и \(\beta\).
    Поскольку \(n\alpha < m\) и \(n > 0\), имеем: \[a < \frac{m}{n}\]
    Кроме того, поскольку \(m \leq n\alpha + 1\) и \(n\alpha + 1 < n\beta\), получаем:
    \[m \leq n\alpha + 1 < n\beta ~ \Rightarrow ~ \frac{m}{n} < \beta\]
    \[\alpha < \frac{m}{n} < \beta\]
    Следовательно, число \(c = \frac{m}{n}\) - искомое.~\qedsymbol
\end{lemma}

\begin{lemma}
    Не существует сечения действительных чисел такого,
    что одновременно в нижнем классе сечения \(A_{*}\) содержится наибольший элемент,
    а в верхнем классе сечения \(A^{*}\) содержится наименьший элемент.

    \vspace{2mm}

    \proofname~ Предположим противное: пусть \(\alpha_{0}\) - наименьшее число в \(A^{*}\), а 
    \(\beta_{0}\) - наибольшее число в \(A_{*}\). Заметим, что по определению сечений 
    \(\alpha_{0} > \beta_{0}\).
    Тогда по Лемме 1.1 найдется такое действительноe число \(c\), что 
    \(\alpha_{0} > c > \beta_{0}\). Но тогда число \(c\) не содержится ни в нижнем,
    ни в верхнем классе сечений - следовательно, не выполняется первое свойство сечений:
    \[A_{*}\cup A^{*} \neq \R\]
    Получили противоречие - следовательно, предположение неверно, и лемма верна.~\qedsymbol
\end{lemma}

\begin{theorem}[Теорема Дедекинда]
    Для любого сечения действительных чисел либо существует наибольший элемент в \(A_{*}\),
    либо существует наименьший элемент в \(A^{*}\). 

    \vspace{2mm}

    \textit{Примечание: теорема не требует доказательства в рамках курса.}
\end{theorem}

\begin{theorem}
    Любому действительному числу взаимно однозначно соответствует бесконечная десятичная дробь.

    \vspace{2mm}

    \textit{Примечание: теорема не требует доказательства в рамках курса.}
\end{theorem}

\Subsection{Теорема о существовании и единственности точной верхней (нижней) 
грани числового множества, ограниченного сверху (снизу)}

\begin{theorem}[Теорема о единственности точной верхней (нижней) грани]
    Если множество \(X\in\R\) имеет точную верхную (нижнюю) грань, то эта грань - единственна.

    \vspace{2mm}

    \proofname~Докажем для точной верхней грани (для случая нижней - аналогично).\\
    Предположим противное: пусть \(\alpha_{0}\) и \(\alpha'_{0}\) - 
    точные верхние грани множества \(X\), причем \(\alpha_{0}\neq\alpha'_{0}\).
    Без ограничения общности будем считать, что \(\alpha_{0}>\alpha'_{0}\).
    Но тогда по определению точной верхней грани (для \(\alpha_{0}\)):
    \[\alpha'_{0} < \alpha_{0} ~ \Rightarrow ~ \exists\,x_{\alpha'_{0}}\in X: 
    \alpha'_{0} < x_{\alpha'_{0}}\]
    Следовательно, \(\alpha'_{0}\) не может являться точной верхней гранью множества \(X\).
    Тогда наше предположение неверно, и точная верхняя грань единственна.~\qedsymbol
\end{theorem}

\begin{theorem}[Теорема о существовании точной верхней (нижней) грани]
    Если множество \(X\in\R\) ограничено сверху (снизу), 
    то оно имеет точную верхную (нижнюю) грань.

    \vspace{2mm}

    \proofname~Докажем для случая ограниченности сверху (для ограниченности снизу - аналогично). \\
    Пусть \(X\in\R\) - ограниченное сверху множество. Рассмотрим 2 случая: \\
        \indent \textbf{1.} Множество \(X\) имеет наибольший элемент. \\
        Тогда этот элемент и будет точной верхней гранью: действительно,
        любой элемент из \(X\) не больше него, и при этом любое число,
        меньшее этого элемента, верхней гранью уже не является.\\
        \indent \textbf{2.} Множество \(X\) не имеет наибольшего элемента.\\
        Построим такие 2 множества: множество \(A^{*}\) будет множеством верхних граней \(X\),
        а множество \(A_{*}\) будет содержать все остальные действительные числа. Тогда:
        \begin{itemize}
            \item \(A_{*} \cup A^{*} = \R\)
            \item \(A_{*} \cap A^{*} = \emptyset\)
            \item \(\forall\,x\in A_{*},~\forall\,y\in A^{*} \longmapsto x < y\)
        \end{itemize}
        Последнее свойство верно потому, что если \(x\geq y\), то \(x\), 
        будучи больше верхней грани, будет сам являться верхней гранью,
        то есть принадлежать \(A^{*}\), чего быть не может.\\
        Тогда по определению наши 2 множества образуют сечение.
        По теореме Дедекинда этому сечению соответствует некоторое число \(\alpha\).
        Заметим, что если какой-то элемент \(x\in X\) принадлежал бы \(A^{*}\),
        то, являясь верхней гранью, он оказался бы наибольшим элементом \(X\),
        чего по нашему предположению быть не может.
        Следовательно, множество \(X\) целиком содержится в нижнем классе сечения \(A_{*}\).\\
        Тогда, если бы \(\alpha\) было наибольшим числом в \(A_{*}\) (по т. Дедекинда),
        то оно было бы верхней гранью \(X\), то есть одновременно принадлежало бы и \(A^{*}\),
        чего быть не может.
        Следовательно, (по т. Дедекинда) \(\alpha\) - наименьшее число в \(A^{*}\).
        Тогда \(\alpha\) - наименьшая из верхних граней - то есть точная верхняя грань.~\qedsymbol
\end{theorem}

\Subsection{Счетность \(\Q\). Несчетность \(\R\)}

\begin{theorem}
    Множество рациональных чисел \(Q\) - счетно.

    \vspace{2mm}

    \proofname~Будем рассматривать только несократимые дроби,
    поскольку сократимые лишь дублируют их. \\
    Рассмотрим положительные дроби вида \(\frac{p}{q}\). Назовем \textit{высотой}
    дроби величину \(H = p+q\) (например, высоте 5 соответствуют дроби
    \(\frac{1}{4}, \frac{2}{3}, \frac{3}{2}, \frac{4}{1}\)). 
    Заметим, что каждой высоте соответствует \textit{конечное} число дробей.
    Расположим дроби следующим образом: по возрастанию высоты, 
    для дробей одной высоты - по возрастсанию величины:
    \[\underbrace{\frac{1}{1}}_{H=1}, \underbrace{\frac{1}{2}, \frac{2}{1}}_{H=2}, \dots\]
    Перенумеруем все дроби, расположенные таким образом:
    номеру \(n\) соответствует дробь \(a_{n}\).
    Для того, чтобы охватить все множество рациональных чисел, зададим 
    следующий ряд чисел:
    \[0, a_{1}, -a_{1}, a_{2}, -a_{2}, \dots \]
    Перенумеруем числа в этом ряду и обозначим как \(b_{n}\).
    Теперь мы взаимно однозначно сопоставили
    каждое натуральное число \(n\) и рациональное число \(b_{n}\).
    Следовательно, по определению \(\Q\) - счетно.~\qedsymbol
\end{theorem}

\begin{lemma}
    Любое бесконечное подмножество счетного множества также счетно.

    \vspace{2mm}

    \proofname~Пусть \(A\) - счетное множество, а множество \(B\subset A\).
    Тогда элементы множества \(A\) можно расположить в порядке возрастания номеров:
    \[A = {a_{1},a_{2},a_{3},a_{4},\dots}\]
    "Пройдемся" по множеству \(A\) в порядке возрастсания номеров.
    Выберем первый элемент, входящий в множество \(B\), и обозначим его \(a_{n_{1}}\).
    Пойдем дальше и аналогично выберем \(a_{n_{2}},a_{n_{3}},\dots\)\\
    Тогда мы взаимно однозначно сопоставили элементы множества \(B\) и натуральные числа:
    \[a_{n_{k}}~\longleftrightarrow~k\]
    Следовательно, множество \(B\) - счетно.
\end{lemma}

\begin{theorem}
    Множество действительных чисел \(\R\) несчетно.

    \vspace{2mm}
    
    По Лемме 1.3 если какое-то подмножество \(\R\) несчетно, то и \(\R\) несчетно.
    Тогда докажем, что отрезок [0, 1] несчетен. \\
    Предположим противное: пусть отрезок [0, 1] счетен. Тогда каждому натуральному номеру \(n\)
    соответствует действительное число, а, следовательно, бесконечная десятичная дробь:
    \[n \longleftrightarrow 0.a_{1n}a_{2n}a_{3n}a_{4n}\dots\]
    Зададим такое действительное число \(\gamma\):
    \[\gamma = 0.b_{1}b_{2}b_{3}b_{4}\dots,\]
    причем \(\forall\,n\in\N\hspace{3mm} b_{n}\neq a_{nn},~ b_{n} \neq 9\).
    Заметим, что \(\gamma\) ни равно ни одному из тех действительных чисел, 
    которым сопоставлен натуральный номер - следовательно, счетность отрезка [0, 1]
    нарушилась, то есть наше предположение неверно, и отрезок [0, 1] несчетен.\\
    Следовательно, несчетно и множество \(\R\).
\end{theorem}
